\section{Diagonal comparison and complementary conditionals}\label{sec:lifting}

The diagonal computation on a Turing object $U$ has two disjoint
domains $a_0,a_1$, one for each output summand. Let $F:\mathscr R\to\mathcal A$
be a range functor. The comparison theorem gives conditions under which
disjoint domains $P\geq F(a_0)$ and $N\geq F(a_1)$ satisfy
$P\vee N\ne1_{FU}$.

For a coproduct $Y+Z$, let
$\varepsilon_0=1_Y+0_Z$ and $\varepsilon_1=0_Y+1_Z$ be its summand
restriction idempotents. For the coproduct injections $i_j$, one has
$\varepsilon_j i_j=i_j$ and $\varepsilon_k i_j=0$ for $k\ne j$.
Thus $\varepsilon_0\varepsilon_1$ vanishes on both injections.
Any upper bound $r$ of $\varepsilon_0,\varepsilon_1$ satisfies
$r i_j=r\varepsilon_j i_j=i_j$ for both $j$. The coproduct property
therefore gives $\varepsilon_0\varepsilon_1=0_{Y+Z}$ and
$\varepsilon_0\vee\varepsilon_1=1_{Y+Z}$.

\subsection{The comparison theorem}

\begin{theorem}[Comparison theorem]\label{thm:comparison}
Let $\mathscr R$ be a recursion category with ranges and binary
restriction coproducts, let $\mathcal A$ be a restriction category with
ranges, restriction zeros, and binary restriction coproducts, and let
$F:\mathscr R\to\mathcal A$ be a range functor.
Choose a Turing morphism $\tau:U\times U\rightharpoonup U$ and an
isomorphism $b:U\to U+U$. Define
\[
 h=b\tau\Delta_U:U\rightharpoonup U+U,\qquad
 a_0=h^*(1_U+0_U),\qquad a_1=h^*(0_U+1_U).
\]
Let $P,N$ be disjoint restriction idempotents on $FU$ such that
$F(a_0)\le P$ and $F(a_1)\le N$. Suppose that
$s:FU\rightharpoonup FU+FU$ satisfies
\[
 s^*(1_{FU}+0_{FU})=N,\qquad
 s^*(0_{FU}+1_{FU})=P,
\]
and that there is a morphism $\ell:U\rightharpoonup U+U$ with
$F\ell=\chi_{U,U}s$.
Then there is a total morphism $g:U\to U$ such that
\[
 (Fg)^*P=(Fg)^*N=0_{FU}.
\]
For any such $g$, if $F(1_U)\ne0$ and inverse image along $Fg$ preserves the join
$P\vee N$, then $P\vee N\ne1_{FU}$.
\end{theorem}
\begin{proof}
Put $u=b^{-1}\ell:U\rightharpoonup U$. Apply the Turing property
to $u\pi_2:U\times U\rightharpoonup U$, obtaining a total morphism
$g:U\to U$ with $\tau(g\times1_U)=u\pi_2$.
Since $g$ is total,
\[
 (g\times1_U)\langle1_U,g\rangle=\langle g,g\rangle=\Delta_Ug,
 \qquad \pi_2\langle1_U,g\rangle=g.
\]
Precomposing with $\langle1_U,g\rangle$ therefore gives
$\tau\Delta_Ug=ug$. Composing with $b$ gives $hg=\ell g$.

Write
$E=Fg:FU\to FU$ and
$H=\chi_{U,U}^{-1}Fh:FU\rightharpoonup FU+FU$.
Applying $F$ to the preceding identity gives $HE=sE$.
For $i=0,1$, let $\varepsilon_i$ and $\varepsilon_i'$ be the $i$th summand
idempotents on $U+U$ and $FU+FU$, respectively. The coproduct comparison
satisfies $F(\varepsilon_i)\chi_{U,U}=\chi_{U,U}\varepsilon_i'$, so
$F(a_i)=H^*\varepsilon_i'$. Pullback along $E$ gives
$E^*F(a_i)=(HE)^*\varepsilon_i'=(sE)^*\varepsilon_i'$.
The two prescribed inverse images under $s$ give
\[
 E^*F(a_0)=E^*N,\qquad E^*F(a_1)=E^*P.
\]
The inequalities imply
\[
 E^*N=E^*F(a_0)\le E^*P,\qquad
 E^*P=E^*F(a_1)\le E^*N.
\]
The two inverse images are equal, and their meet is
$E^*(PN)=0_{FU}$. Hence both are zero.
If $P\vee N=1_{FU}$, totality of $E$ and preservation of this join
would give
\[
 1_{FU}=E^*1_{FU}=E^*(P\vee N)
       =E^*P\vee E^*N=0_{FU},
\]
contrary to $F(1_U)\ne0$.
\end{proof}

\begin{corollary}[Lifted separators]\label{cor:separator}
Under the hypotheses of Theorem~\ref{thm:comparison} on
$\mathscr R$, $\mathcal A$, $F$, $\tau$, and $b$, suppose that
$F(1_U)\ne0$ and inverse images along total morphisms in $\mathcal A$ preserve
binary joins. A complementary pair $P,N$ on $FU$ satisfying
$F(a_0)\le P$ and $F(a_1)\le N$ admits no pair of morphisms
$s:FU\rightharpoonup FU+FU$ and
$\ell:U\rightharpoonup U+U$ with
\[
 s^*(1_{FU}+0_{FU})=N,\qquad
 s^*(0_{FU}+1_{FU})=P,\qquad
 F\ell=\chi_{U,U}s.
\]
\end{corollary}
\begin{proof}
Such morphisms would give $P\vee N\ne1_{FU}$ by
Theorem~\ref{thm:comparison}, contradicting complementarity.
\end{proof}

\subsection{Complementary conditionals}

\begin{corollary}[Inseparability from complementary conditionals]
\label{cor:source-conditionals}
Under the hypotheses of Theorem~\ref{thm:comparison} on
$\mathscr R$, $\mathcal A$, $F$, $\tau$, and $b$, suppose that
$F(1_U)\ne0$ and inverse images along total morphisms in $\mathcal A$ preserve
binary joins of restriction idempotents.
Suppose that, for every complementary pair $c,c'$ on $U$, there
is a total morphism $\ell_c:U\to U+U$ satisfying
\[
 \ell_cc=i_1c,\qquad \ell_cc'=i_0c',
\]
where $i_0,i_1:U\to U+U$ are the coproduct injections.
Then the restriction idempotents $a_0,a_1$ of
Theorem~\ref{thm:comparison} are inseparable in $\mathscr R$.
\end{corollary}
\begin{proof}
Put $\varepsilon_0=1_U+0_U$ and $\varepsilon_1=0_U+1_U$.
The diagonal domains are disjoint, since
\[
 a_0a_1=h^*(\varepsilon_0\varepsilon_1)=h^*0_{U+U}=0_U.
\]
Suppose $c,c'$ separate $a_0,a_1$. The defining equations of
$\ell_c$ give
\[
 (\varepsilon_0\ell_c)c=0=(i_0c')c,\qquad
 (\varepsilon_0\ell_c)c'=i_0c'=(i_0c')c'.
\]
Thus $\varepsilon_0\ell_c$ and $i_0c'$ agree after precomposition
with $[c,c']:U+U\rightharpoonup U$. Its range is
$c\vee c'=1_U$, so the agreement law gives $\varepsilon_0\ell_c=i_0c'$.
Interchanging the summands gives $\varepsilon_1\ell_c=i_1c$.
Taking restrictions yields $\ell_c^*\varepsilon_0=c'$ and
$\ell_c^*\varepsilon_1=c$.

Since $F$ preserves meets, joins, zeros, and identities,
$P=F(c)$ and $N=F(c')$ are complementary on $FU$.
For the canonical coproduct isomorphism
$\chi_{U,U}:FU+FU\to F(U+U)$, define
$s_c=\chi_{U,U}^{-1}F(\ell_c):FU\to FU+FU$.
Restriction preservation and the coproduct comparison give
\[
 s_c^*(1_{FU}+0_{FU})=N,\qquad
 s_c^*(0_{FU}+1_{FU})=P.
\]
Together with $F(\ell_c)=\chi_{U,U}s_c$, these identities contradict
Corollary~\ref{cor:separator}.
\end{proof}

\begin{proposition}\label{prop:disjoint-joins}
Let $\mathcal C$ be a restriction category with restriction zeros and
binary restriction coproducts. Suppose every pair of parallel morphisms
whose restriction idempotents are disjoint has a join in the restriction
order.
For complementary restriction idempotents $c,c'$ on $X$, the morphism
\[
 \ell_c=(i_1c)\vee(i_0c'):X\rightharpoonup X+X
\]
is total and satisfies $\ell_cc=i_1c$ and $\ell_cc'=i_0c'$.
\end{proposition}
\begin{proof}
The restrictions of $i_1c$ and $i_0c'$ are $c$ and $c'$ because
the injections are total. The upper-bound equations for the join
give $\ell_cc=i_1c$ and $\ell_cc'=i_0c'$.
Both $c$ and $c'$ lie below $\bar\ell_c$, so
$1_X=c\vee c'\le\bar\ell_c$ and $\ell_c$ is total.
\end{proof}

Appendix~\ref{app:lifting-counterexample} gives a faithful range functor
and a complementary target separator whose conditional has no source lift.

\subsection{Di Paola and Heller's diagonal theorem}
\label{subsec:di-paola-heller-diagonal}

Di Paola and Heller assume throughout their Section~8 that the category
is a recursion category \cite[p.~621]{DPH}. For Theorem~8.15 they
assume a dominical category with binary coproducts and, for each
$f:Y\rightharpoonup Z$, a section $s:Z\rightharpoonup Y$ satisfying
$fs=\bar s$ and $fsf=f$. The restriction idempotent $fs=\bar s$
fixes $f$, and $uf=vf$ implies $u(fs)=v(fs)$ for parallel
$u,v:Z\rightharpoonup W$. Thus $\ran(f)=fs$
\cite[p.~612, Definitions~6.1--6.2]{DPH}.
They also assume that $f+g$ is total whenever $f$ and $g$ are total
and that the canonical morphisms
$(X\times Y)+(X\times Z)\to X\times(Y+Z)$ are isomorphisms
\cite[p.~615, Definition~7.1]{DPH}.

Let $\tau:X\times X\rightharpoonup X+X$ be a Turing morphism in their
sense: for every $p:X\times X\rightharpoonup X+X$ there is a total
$k:X\to X$ with $p=\tau(k\times 1_X)$ \cite[p.~601, Definition~4.1]{DPH}.
Put $h=\tau\Delta_X$. Its two summand inverse images
$h^*(1_X+0_X)$ and $h^*(0_X+1_X)$ are disjoint. If a complementary
pair $c,d$ separated them, their diagonal argument would give a total
$g:X\to X$ with $g^*c=g^*d=0_X$
\cite[pp.~625--626, Theorem~8.15]{DPH}. Then $cg=dg=0_X$, and range
agreement gives $c\ran(g)=d\ran(g)=0_X$. Since restriction idempotents
commute, $\ran(g)[c,d]=0$ for the copairing $[c,d]:X+X\rightharpoonup X$.
Now $\ran[c,d]=c\vee d=1_X$, so range agreement gives $\ran(g)=0_X$.
Hence $g=0_X$, and totality gives $1_X=\bar g=0_X$. The two diagonal
domains are therefore inseparable when $1_X\ne0_X$.
In Theorem~\ref{thm:comparison}, the lift of the conditional makes the
inverse images of $P,N$ along $Fg$ vanish. If $F(1_U)\ne 0$ and
$(Fg)^*$ preserves $P\vee N$, this gives $P\vee N\ne 1_{FU}$.
